\section{Regular windows, pair profiles and compression}
\label{p1:arith:section}

The first correction to the Poisson law comes from two windows whose integer
arguments follow the same rational progression.  The relevant quantity is the
number of square relations, with their affine signs retained.  Counting only
the pairs that support a relation would lose both the size of the correction
and its dependence on the integer length.

Throughout this section, $I_N=[N,2N)\cap\mathbb Z$, $B=L+1$,
$p=2^{-L}$ and $\lambda=Np$, with $|L-\log_2N|\leq C$.
All constants are uniform in this band.  The square-class vector of a positive
integer $n$ is denoted by $v(n)$; its prime coordinates are $v_\ell(n)\bmod2$.
For $V_x=\{x-1,\ldots,x+L-1\}$, let $E_x$ be the space of even subsets of
$V_x$, and let $S_x$ be its image under the square-class map.  Thus $E_x$ has
dimension $L$.  A selection $A\in E_x$ has affine sign
$(-1)^{\ind_{\{x-1\in A\}}}$ for the run-start event $J_x=1$.

For the arithmetic profiles in this section, write
\[
 \rho(I)=|I|L-\dim\sum_{x\in I}S_x,\qquad
 q(I)=\P(J_x=1\text{ for every }x\in I).
\]
Here $\rho$ includes the local kernels. On the regular set it equals
$\rho_{\mathrm{rel}}$; on arbitrary windows the two quantities must be
distinguished. The exact-marginal probability argument uses
$\rho_{\mathrm{rel}}$.
The affine character on the kernel is either trivial or nontrivial.  Accordingly,
with $\eta(I)\in\{0,1\}$ denoting compatibility, finite Fourier inversion gives
\begin{equation}\label{p1:arith:affine}
 q(I)=p^{|I|}\eta(I)2^{\rho(I)}.
\end{equation}
This formula includes incompatible events.  In particular,
$|\eta(I)2^{\rho(I)}-1|\leq2^{\rho(I)}-1$ whenever the individual systems
have full rank.  For a pair, the same inequality holds even before this
restriction, since a zero-dimensional kernel has no compatibility obstruction.

\subsection{The regular set and the cost of its complement}

For $H\geq2$ put
\[
 K_H(n)=\prod_{\substack{\ell>H\ {\rm prime}\\v_\ell(n)\ {\rm odd}}}\ell.
\]
An integer with $K_H(n)=1$ will be called $H$-defective.  When the diameter of
an interval is less than $H$, every other integer in it has a prime coordinate
which occurs nowhere else in the interval.  Such a coordinate excludes that
integer from every square relation.

\begin{proposition}[Regular set]\label{p1:arith:core}
Let
\[
 G_N=\{x\in I_N:\dim S_x=L\},\qquad Z_N^\circ=\sum_{x\in G_N}J_x.
\]
Then $|I_N\setminus G_N|=N^{o(1)}$, and
\begin{equation}\label{p1:arith:coretransfer}
 d_{\rm TV}(\Law(Z_N),\Law(Z_N^\circ))+
 |\E Z_N-\lambda|+|\E Z_N^\circ-\lambda|
 \ll_C N^{-1+o(1)}.
\end{equation}
More generally, for every deterministic $A\subset I_N$ and every choice of
affine targets on its windows,
\begin{equation}\label{p1:arith:marginalsum}
 \sum_{x\in A}|\P(A_xX=b_x)-p|\leq pN^{o(1)}.
\end{equation}
If the fixed moments of $Z_N$ are uniformly bounded, then, for every fixed
integer $j\geq1$ and every $\varepsilon>0$,
\begin{equation}\label{p1:arith:momenttransfer}
 \big|\E Z_N^j-\E(Z_N^\circ)^j\big|
 \ll_{C,j,\varepsilon}N^{-1+\varepsilon}.
\end{equation}
In particular the second factorial cumulants have the same transfer bound.
\end{proposition}

The proof is given in Section~\ref{p14:comp:regular-transfer}.

\subsection{Sparse hosts and weighted nonaligned pairs}

For separated windows, $|x-y|>L$, a prime larger than $B$ occurs with odd
valuation at most once in either window.  Its equation either pins one
vertex to zero or joins two vertices.  After the pinned components are
removed, each remaining component has one binary degree of freedom.  If
$D(x,y)$ counts the defective vertices and $c(x,y)$ the other unpinned
components, then
\begin{equation}\label{p1:arith:graphbound}
 \rho(x,y)\leq D(x,y)+c(x,y),\qquad D(x,y)=o(B).
\end{equation}
The small prime equations and the two block parities can only reduce this
dimension.

A pair is called aligned if, in one orientation, it satisfies
$bx-ay=h$ for coprime positive $a,b\leq B^3$ and
$|h|\leq(a+b)B$.  Using this slightly enlarged definition includes every
alignment produced by two exact units.  Two different reduced slopes cannot
occur for the same dyadic pair once $N$ is large: distinct slopes are separated
by at least $B^{-6}$, whereas the allowed displacement changes the ratio of
the starts by $O(B^{O(1)}/N)$.

\begin{proposition}[Support and residual mass]\label{p1:arith:hosts}
For every $\varepsilon>0$,
\begin{align}
 \#\{(x,y)\in I_N^2:|x-y|>L,\ \rho(x,y)>0\}
   &\ll_{C,\varepsilon}N^{1+\varepsilon},
       \label{p1:arith:hostbound}\\
 \sum_{\substack{x,y\in I_N,\ |x-y|>L\\(x,y)\ {\rm nonaligned}}}
       (2^{\rho(x,y)}-1)
   &\ll_{C,\varepsilon}N^{1+\varepsilon}.
       \label{p1:arith:nonaligned}
\end{align}
\end{proposition}

The elementary proof in Section~\ref{p14:comp:sparse-hosts} first gives
the sufficient bound $N^{6/5+\varepsilon}$ in both displays.
The stronger bounds stated here are completed in
Section~\ref{p14:comp:confinement}, using the conditional half-count
and dense-pair confinement. That later strengthening is not an input
to either of those arguments; the elementary bounds suffice for the
earlier pair-profile calculation.

\subsection{Good edges, compression, and one additional window}

Let $\mathcal E_N$ consist of the unordered compatible dependent pairs in
the channels $2:3$ and $3:4$ for which the distinct normalized values have private
prime coordinates above the fixed threshold $H$.  Write
$q_e=\P(J_x=J_y=1)$ for $e=\{x,y\}$ and
$\beta_N=\sum_{e\in\mathcal E_N}q_e$.  The family is deterministic.

\begin{corollary}[The pair input for the expansion]\label{p1:arith:edgefamily}
One has $|\mathcal E_N|=O_C(NB)$ and, for an edge in channel $a:b$,
\[
 \rho(e)\leq L/b+O(1),\qquad q_e\leq p^2 2^{L/b+O(1)}.
\]
Moreover,
\begin{align}
 2\beta_N&=\alpha_N+O_{C,\varepsilon}(N^{-4/5+\varepsilon}),
    \label{p1:arith:edgemass}\\
 \sum_{\substack{x,y\in G_N,\ x\ne y\\\{x,y\}\notin\mathcal E_N}}
   |\E[(J_x-p)(J_y-p)]|
 &\ll_{C,\varepsilon}N^{-4/5+\varepsilon}+B/N.
    \label{p1:arith:badpairs}
\end{align}
In particular $\beta_N\asymp_C LN^{-2/3}$.
\end{corollary}

The proof is given in Section~\ref{p14:comp:good-edge-mass}.

\begin{lemma}[Exact compression of a good edge]\label{p1:arith:compression}
For a good edge in channel $a:b$, the joint homogeneous event space is the
injective square-class image of the even-selection space of
\[
 U_e=bV_x\cup aV_y,
\]
where repeated values are included only once.  In particular its rank is
\begin{equation}\label{p1:arith:compressedrank}
 \dim(S_x+S_y)=|U_e|-1\geq(2-1/b)L-O(1).
\end{equation}
After the endpoint offsets are fixed, $U_e$ is a subset of an interval of
length $O(B)$ depending on one integer translation parameter.  There are
only $B^{O(1)}$ offset patterns.
\end{lemma}

\begin{proof}
Multiplying every integer in a window by the same constant does not change
the square class of an even selection.  The two even-selection spaces on
$bV_x$ and $aV_y$ therefore map to the original event spaces.  They intersect
on at least one normalized value: dependence in the good family requires
at least two exact units.  The sum of the even-selection spaces of two
intersecting finite sets is the even-selection space of their union.  Its
square-class map is injective, because every distinct value in $U_e$ has a
private prime.  There are at most $L/b+O(1)$ repeated values, which proves
\eqref{p1:arith:compressedrank}.
\end{proof}


\subsection{Pair profiles and the regular-set estimates}
\label{p14:comp:pair-profiles}

This section supplies the estimates used in Section~\ref{p1:arith:section}.
The regular-set transfer is used for the explicit channel formulas and for
the independent fixed-moment argument. The scalar reconstruction with exact marginals
does not discard these exceptional windows.

\subsubsection{Regular-set transfer}
\label{p14:comp:regular-transfer}

\begin{proof}[Proof of Proposition~\ref{p1:arith:core}]
The one-window estimates proved in Lemma~\ref{p1:arith:defects} give
$m_B(x)=O_C(B/\log B)$ uniformly, and only $N^{o(1)}$ windows contain two
$B$-defective integers.  A nonzero relative relation is an even nonempty
selection.  Private prime coordinates therefore force it to use at least two
defects.  Its dimension is at most $m_B(x)$.  Consequently
\[
 \sum_{x\in I_N}(2^{L-\dim S_x}-1)=N^{o(1)}.
\]
Fourier inversion proves \eqref{p1:arith:marginalsum}.  The sum of the
probabilities of active discarded starts is at most $pN^{o(1)}$; the coupling
using the original prime signs proves the total-variation assertion.  On $G_N$
the marginal is exactly $p$, giving both assertions about the mean.

For the last statement, let $D=\{Z_N\ne Z_N^\circ\}$.  Since
$0\leq Z_N^\circ\leq Z_N$, H\"older's inequality gives, for any fixed $s>1$,
\[
 \E\big[Z_N^j\ind_D\big]
 \leq(\E Z_N^{js})^{1/s}\P(D)^{1-1/s}.
\]
Choose a sufficiently large integer $s$, and then use the fixed-moment bound.
This last step is used only after that bound has been established; a small
coupling error alone would not justify a variance transfer.
\end{proof}

\subsubsection{Sparse hosts and weighted residual mass}
\label{p14:comp:sparse-hosts}

\begin{proof}[The elementary bounds preceding Proposition~\ref{p1:arith:hosts}]
A relative relation that is internal to one window contributes at most
$N^{1+o(1)}$ hosts, by Proposition~\ref{p1:arith:core}.  Otherwise it selects
at least two vertices in each window.  Put $T=N^{3/5}$.  If all selected
kernels are at most $T$, each window contains two small kernels.  The estimate
\begin{equation}\label{p1:arith:twokernelused}
 \#\{x\in I_N:\#\{n\in V_x:K_B(n)\leq N^\beta\}\geq2\}
   \ll_{C,\beta,\varepsilon}N^{\max(\beta,7/12)+\varepsilon}
\end{equation}
proved in Lemma~\ref{p1:arith:twokernel} gives $N^{6/5+\varepsilon}$
choices.  If one selected kernel exceeds $T$, Lemma~\ref{p1:arith:anchor}
gives
$N^{3/2+o(1)}T^{-1/2+\tau}$, for every fixed $\tau>0$.
Choose $\tau$ sufficiently small.

For the weighted assertion, first restrict to $c(x,y)<\eta B$.
Equation~\eqref{p1:arith:graphbound} bounds the weight by
$2^{\eta B+o(B)}$; choose $\eta$ after the desired error exponent.
The complementary sector has total weight $O_{C,\eta,\varepsilon}(N^{1+\varepsilon})$
by Lemma~\ref{p1:arith:dense}.  The latter estimate uses the component
equations themselves, rather than just their support: many small components
put the pair on a multiquadratic surface, and its low-degree curves are either
aligned or absent from the dyadic region.
This proves both bounds with exponent $6/5+\varepsilon$, which is all
that the pair-profile proof below uses. The exponent $1+\varepsilon$
in the revised statement is proved later in
Section~\ref{p14:comp:confinement}; no late confinement estimate is
being inferred from this elementary calculation.
\end{proof}

\subsubsection{The complete pair profile}
\label{p14:comp:complete-profile}

\begin{proof}[Proof of Theorem~\ref{p1:arith:pairs}]
For the channel $2:3$, either orientation permits
$N/2\leq t<2N/3$, up to $O(B)$ endpoint changes.  For $3:4$ the corresponding
interval is $N/3\leq t<N/2$.  Both intervals have length $N/6+O(B)$.
The two orientations therefore give $N/3$ times the finite sums~\eqref{p1:arith:S23}--\eqref{p1:arith:S34}.
Translations near a boundary contribute $O(B^3 2^{L/b})$.

For either fixed channel choose $H$ a sufficiently large fixed multiple of
$B$ and normalize the values to $bV_x\cup aV_y$.  Its diameter is $O(B)$.
If its distinct values are all nondefective at level $H$, each has a private
prime and the only relations before normalization are the even selections
of exact units.  Lemma~\ref{p1:arith:defects} shows that the exceptional
channel pairs number $N^{1/2+o(1)}$.  Lemma~\ref{p1:arith:alignedrank} bounds
their nullity by $L/b+o(L)$.  Replacing the rational kernel by the full
arithmetic kernel consequently costs at most
$N^{1/2+1/b+o(1)}$, for either the positive or the signed sum.

There are $N B^{O(1)}$ aligned hosts altogether.  The slope of height two
forces both starts within $O(B)$ of the dyadic endpoints; its mass is
$N^{1/2+o(1)}$.  Height one has no exact unit in separated windows.
The only reduced slopes of heights three and four that occur in the interior
are $2:3$ and $3:4$.  Every remaining channel has height at least five and
nullity at most $L/5+o(L)$, again by Lemma~\ref{p1:arith:alignedrank}; its
total contribution is $N^{6/5+o(1)}$. The elementary version just
proved controls the nonaligned sector at the same sufficient scale.
The inequality following
\eqref{p1:arith:affine} permits exactly the same error estimates for the signed
sum.

At distances $1,\ldots,L-1$, two run starts are incompatible.  At distance
$L$, their edge systems form a tree on an interval of diameter $2L$;
Lemma~\ref{p1:arith:defects}, at threshold $2L+1$, bounds its nullity by
$o(L)$.  Hence the entire contact band changes the second factorial moment
by $O(N^{-1+o(1)})$.  Multiplying \eqref{p1:arith:signedprofile} by $p^2$,
and using $\mu_N-\lambda=O(N^{-1+o(1)})$, proves
\eqref{p1:arith:cumulantexpansion}.  Finally
$2^{(L-r)/b}=2^{-r/b}(N/\lambda)^{1/b}$ in
\eqref{p1:arith:S23}--\eqref{p1:arith:S34} gives the two equivalents.
\end{proof}

\subsubsection{The good-edge mass}
\label{p14:comp:good-edge-mass}

\begin{proof}[Proof of Corollary~\ref{p1:arith:edgefamily}]
Each start has $O(B)$ possible partners in a fixed channel.  The rational
kernel has dimension one less than the number of exact units when that
number is positive, giving the rank bound.  All pairs excluded from the good
family have precisely the error mass estimated in the proof of
Theorem~\ref{p1:arith:pairs}, together with the contact band.
For good compatible pairs the centered contribution is $q_e-p^2$.
Since $p^2|\mathcal E_N|=O(B/N)$, replacing it by $q_e$ has the stated cost.
The asymptotic follows from \eqref{p1:arith:leadingphase}.
\end{proof}

